        %%******************************************%%
        %%                                          %%
        %%        Modello di tesi di laurea         %%
        %%            di Andrea Giraldin            %%
        %%                                          %%
        %%             2 novembre 2012              %%
        %%                                          %%
        %%******************************************%%


% I seguenti commenti speciali impostano:
% 1.
% 2. PDFLaTeX come motore di composizione;
% 3. tesi.tex come documento principale;
% 4. il controllo ortografico italiano per l'editor.

% !TEX encoding = UTF-8
% !TEX TS-program = pdflatex
% !TEX root = tesi.tex
% !TEX spellcheck = it-IT

\documentclass[10pt,                    % corpo del font principale
               a4paper,                 % carta A4
               twoside,                 % impagina per fronte-retro
               openright,               % inizio capitoli a destra
               english,
               italian,
               ]{book}

%**************************************************************
% Importazione package
%**************************************************************

%\usepackage{amsmath,amssymb,amsthm}    % matematica

\usepackage[T1]{fontenc}                % codifica dei font:
                                        % NOTA BENE! richiede una distribuzione *completa* di LaTeX

\usepackage[utf8]{inputenc}             % codifica di input; anche [latin1] va bene
                                        % NOTA BENE! va accordata con le preferenze dell'editor

\usepackage[english, italian]{babel}    % per scrivere in italiano e in inglese;
                                        % l'ultima lingua (l'italiano) risulta predefinita

\usepackage{bookmark}                   % segnalibri

\usepackage{caption}                    % didascalie

\usepackage{chngpage,calc}              % centra il frontespizio

\usepackage{csquotes}                   % gestisce automaticamente i caratteri (")

\usepackage{emptypage}                  % pagine vuote senza testatina e piede di pagina

\usepackage{epigraph}			% per epigrafi

\usepackage{eurosym}                    % simbolo dell'euro

%\usepackage{indentfirst}               % rientra il primo paragrafo di ogni sezione

\usepackage{graphicx}                   % immagini

\usepackage{hyperref}                   % collegamenti ipertestuali

\usepackage[binding=5mm]{layaureo}      % margini ottimizzati per l'A4; rilegatura di 5 mm

\usepackage{listings}                   % codici

\usepackage{microtype}                  % microtipografia

\usepackage{mparhack,fixltx2e,relsize}  % finezze tipografiche

\usepackage{nameref}                    % visualizza nome dei riferimenti

\usepackage[font=small]{quoting}        % citazioni

\usepackage{subfig}                     % sottofigure, sottotabelle

\usepackage[italian]{varioref}          % riferimenti completi della pagina

\usepackage[dvipsnames]{xcolor}         % colori

\usepackage{booktabs}                   % tabelle
\usepackage{tabularx}                   % tabelle di larghezza prefissata
\usepackage{longtable}                  % tabelle su più pagine
\usepackage{ltxtable}                   % tabelle su più pagine e adattabili in larghezza

\usepackage[toc, acronym]{glossaries}   % glossario
                                        % per includerlo nel documento bisogna:
                                        % 1. compilare una prima volta tesi.tex;
                                        % 2. eseguire: makeindex -s tesi.ist -t tesi.glg -o tesi.gls tesi.glo
                                        % 3. eseguire: makeindex -s tesi.ist -t tesi.alg -o tesi.acr tesi.acn
                                        % 4. compilare due volte tesi.tex.

\usepackage[backend=biber,style=verbose-ibid,hyperref,backref]{biblatex}
                                        % eccellente pacchetto per la bibliografia;
                                        % produce uno stile di citazione autore-anno;
                                        % lo stile "numeric-comp" produce riferimenti numerici
                                        % per includerlo nel documento bisogna:
                                        % 1. compilare una prima volta tesi.tex;
                                        % 2. eseguire: biber tesi
                                        % 3. compilare ancora tesi.tex.

%**************************************************************
% file contenente le impostazioni della tesi
%**************************************************************

%**************************************************************
% Frontespizio
%**************************************************************

% Autore
\newcommand{\myName}{Timoty Granziero}
\newcommand{\myTitle}{Approccio a Microservizi nello sviluppo di Web Application moderne​}

% Tipo di tesi
\newcommand{\myDegree}{Tesi di laurea triennale}

% Università
\newcommand{\myUni}{Università degli Studi di Padova}

% Facoltà
\newcommand{\myFaculty}{Corso di Laurea in Informatica}

% Dipartimento
\newcommand{\myDepartment}{Dipartimento di Matematica "Tullio Levi-Civita"}

% Titolo del relatore
\newcommand{\profTitle}{Prof.}

% Relatore
\newcommand{\myProf}{Gilberto File'}

% Luogo
\newcommand{\myLocation}{Padova}

% Anno accademico
\newcommand{\myAA}{2018-2019}

% Data discussione
\newcommand{\myTime}{18 Luglio 2019}


%**************************************************************
% Impostazioni di impaginazione
% see: http://wwwcdf.pd.infn.it/AppuntiLinux/a2547.htm
%**************************************************************

\setlength{\parindent}{14pt}   % larghezza rientro della prima riga
\setlength{\parskip}{0pt}   % distanza tra i paragrafi


%**************************************************************
% Impostazioni di biblatex
%**************************************************************

\bibliography{bibliografia} % database di biblatex

\defbibheading{bibliography} {
    \cleardoublepage
    \phantomsection
    \addcontentsline{toc}{chapter}{\bibname}
    \chapter*{\bibname\markboth{\bibname}{\bibname}}
}

\setlength\bibitemsep{1.5\itemsep} % spazio tra entry

\DeclareBibliographyCategory{opere}
\DeclareBibliographyCategory{web}

\addtocategory{opere}{womak:lean-thinking}
\addtocategory{web}{site:agile-manifesto}

\defbibheading{opere}{\section*{Riferimenti bibliografici}}
\defbibheading{web}{\section*{Siti Web consultati}}


%**************************************************************
% Impostazioni di caption
%**************************************************************
\captionsetup{
    tableposition=top,
    figureposition=bottom,
    font=small,
    format=hang,
    labelfont=bf
}

%**************************************************************
% Impostazioni di glossaries
%**************************************************************

%**************************************************************
% Acronimi
%**************************************************************
%\renewcommand{\acronymname}{Acronimi e abbreviazioni}

\newacronym[description={\glslink{httpg}{HyperText Transfer Protocol}}]
    {http}{HTTP}{HyperText Transfer Protocol}

\newacronym[description={\glslink{apig}{Application Program Interface}}]
    {api}{API}{Application Program Interface}

\newacronym[description={\glslink{umlg}{Unified Modeling Language}}]
    {uml}{UML}{Unified Modeling Language}

\newacronym[description={\glslink{ictg}{Information and Communications Technology}}]
    {ict}{ICT}{Information and Communications Technology}

\newacronym[description={\glslink{isog}{International Organization for Standardization}}]
    {iso}{ISO}{International Organization for Standardization}

\newacronym[description={\glslink{ohsasg}{Occupational Health and Safety Assessment Series}}]
    {ohsas}{OHSAS}{Occupational Health and Safety Assessment Series}

\newacronym{bss}{BSS}{Supporto al Business}

\newacronym{soa}{SOA}{Società Organismi di Attestazione}

\newacronym{eai}{EAI}{Enterprise Application Integration}

\newacronym{it}{IT}{Information Technology}

\newacronym{mit}{MIT}{Massachusetts Institute of Tecnology}

\newacronym[description={\glslink{iotg}{Internet of Things}}]
    {iot}{IoT}{Internet of Things}

\newacronym[description={\glslink{jvmg}{Java Virtual Machine}}]
    {jvm}{JVM}{Java Virtual Machine}

\newacronym[description={\glslink{dbmsg}{Database Management System}}]
    {dbms}{DBMS}{Database Management System}

\newacronym[description={\glslink{jsong}{JavaScript Object Notation}}]
    {json}{JSON}{JavaScript Object Notation}

\newacronym{smtp}{SMTP}{Simple Mail Transfer Protocol}

\newacronym[description={\glslink{restg}{Representational State Transfer}}]
    {rest}{REST}{Representational State Transfer}

\newacronym{ieee}{IEEE}{Institute of Electrical and Electronics Engineers}

\newacronym{ui}{UI}{User Interface}

\newacronym{ide}{IDE}{Integrated Development Environment}

\newacronym{ux}{UX}{User Experience}

\newacronym{html}{HTML}{HyperText Markup Language}

\newacronym{css}{CSS}{Cascading Style Sheets}

\newacronym{ram}{RAM}{Random Access Memory}

\newacronym{cli}{CLI}{Command Line Interface}

\newacronym{vcs}{VCS}{Version Control System}

\newacronym{dql}{DQL}{Database Query Language}

\newacronym{jdk}{JDK}{Java Development Kit}

\newacronym{cfu}{CFU}{Credito Formativo Universitario}

\newacronym{ioc}{IoC}{Inversion of Control}

%**************************************************************
% Glossario
%**************************************************************

%\renewcommand{\glossaryname}{Glossario}

\newglossaryentry{apig}
{
    name=API,
    text=Application Program Interface,
    sort=api,
    description={In informatica con il termine \emph{Application Programming Interface API} (ing. interfaccia di programmazione di un'applicazione) si indica ogni insieme di procedure disponibili al programmatore, di solito raggruppate a formare un set di strumenti specifici per l'espletamento di un determinato compito all'interno di un certo programma. La finalità è ottenere un'astrazione, di solito tra l'hardware e il programmatore o tra software a basso e quello ad alto livello semplificando così il lavoro di programmazione}
}

\newglossaryentry{umlg}
{
    name=UML,
    text=UML,
    sort=uml,
    description={In ingegneria del software \emph{UML, Unified Modeling Language} (ing. linguaggio di modellazione unificato) è un linguaggio di modellazione e specifica basato sul paradigma object-oriented. L'\emph{UML} svolge un'importantissima funzione di ``lingua franca'' nella comunità della progettazione e programmazione a oggetti. Gran parte della letteratura di settore usa tale linguaggio per descrivere soluzioni analitiche e progettuali in modo sintetico e comprensibile a un vasto pubblico}
}

\newglossaryentry{ictg}
{
    name=Information and Communications Technology,
    text=ICT,
    sort=ict,
    description={in informatica, il termine \emph{Information and Communications Technology} sono l'insieme dei metodi e delle tecniche utilizzate nella trasmissione, ricezione ed elaborazione di dati e informazioni}
}

\newglossaryentry{isog}
{
    name=ISO,
    text=ISO,
    sort=iso,
    description={L'\emph{Organizzazione internazionale per la normazione} è l'organizzazione più importante a livello mondiale per la definizione di \textit{standard}}
}

\newglossaryentry{ohsasg}
{
    name=Occupational Health and Safety Assessment Series,
    text=OHSAS,
    sort=ohsas,
    description={L'\emph{Occupational Health and Safety Assessment Series} identifica uno standard inglese per un sistema di gestione e sicurezza dei lavoratori}
}

\newglossaryentry{big-data}
{
    name=Big Data,
    text=Big Data,
    sort=big data,
    description={Il termine si riferisce a una raccolta di dati (strutturati e non) talmente estesa in termini di velocità, volume e varietà da richiedere apposite tecnologie e metodi per l'estrazione. Ciò che conta non è la quantità di dati, ma come essi vengono utilizzati}
}

\newglossaryentry{iotg}
{
    name=Internet of Things,
    text=IOT,
    sort=iot,
    description={Nelle telecomunicazioni, il termine si riferisce all'estensione di internet al mondo delle cose e luoghi concreti}
}

\newglossaryentry{webapp}
{
    name=Web Application,
    text=web application,
    sort=web application,
    description={Un applicazione web è un applicazione distribuita via web, pertanto accessibile per mezzo della rete e di un \emph{browser}}
}

\newglossaryentry{skill-matrix}
{
    name=Skill Matrix,
    text=skill matrix,
    sort=skill matrix,
    description={Una matrice (e.g. un foglio di calcolo) usata dall'azienda Sync Lab per documentare le competenze dei candidati interessati all'assunzione, in cui è riportato il nome della competenza, il livello e la categoria. Il documento viene compilato dai candidati prima di un eventuale colloquio}
}

\newglossaryentry{angular}
{
    name=Angular,
    text=Angular,
    sort=angular,
    description={Framework \textit{open source} per lo sviluppo di applicazioni web, rilasciato con licenza \gls{mit}, nato come evoluzione di AngularJS. Il linguaggio di programmazione usato da Angular è \gls{typescript}, a differenza del predecessore che usava \gls{javascript}}
}

\newglossaryentry{spring}
{
    name=Spring,
    text=Spring,
    sort=spring,
    description={Framework \textit{open source} per lo sviluppo di applicazioni con target principale Java, ma supporta ufficialmente anche Kotlin e Groovy. A Spring sono associati molti sotto-progetti, come Spring Boot e Spring Cloud, sviluppati per fornire modularità al framwork}
}

\newglossaryentry{java}
{
    name=Java,
    text=Java,
    sort=java,
    description={Java è un linguaggio di programmazione ad oggetti e a tipizzazione statica \emph{general purpose} nato a metà degli anni 90 che si appoggia alla \gls{jvm}}
}

\newglossaryentry{jvmg}
{
    name=Java Virtual Machine,
    text=Java Virtual Machine,
    sort=jvm,
    description={Componente della piattaforma \gls{java} che esegue programmi tradotti in \emph{bytecode} dopo una fase di compilazione}
}

\newglossaryentry{dbmsg}
{
    name=Database Management System,
    text=DBMS,
    sort=dbms,
    description={In informatica, il termine \gls{dbms} indica un sistema software progettato per consentire la manipolazione, la creazione e l'interrogazione efficiente su un database}
}

\newglossaryentry{mongodb}
{
    name=MongoDB,
    text=MongoDB,
    sort=MongoDB,
    description={È un \gls{dbms} non relazionale (NoSQL) orientato ai documenti: ciò significa che abbandona la struttura tradizionale basata su tabelle dei database relazionali in favore di documenti in stile \gls{json}}
}

\newglossaryentry{jsong}
{
    name=JSON,
    text=JSON,
    sort=JSON,
    description={Acronimo di JavaScript Object Notation, è un formato usato nel web adatto allo scambio dei dati fra applicazioni client-server. Si basa sul linguaggio \gls{javascript} da cui prende spunto per la sintassi}
}

\newglossaryentry{cloud-computing}
{
    name=Cloud Computing,
    text=Cloud Computing,
    sort=Cloud Computing,
    description={È un paradigma di erogazione di servizi \emph{on demand} offerti ad un cliente da un fornitore attraverso il web, a partire da
    un sistema con risorse configurabili e preesistenti, generalmente disponibili in remoto}
}

\newglossaryentry{javascript}
{
    name=JavaScript,
    text=JavaScript,
    sort=JavaScript,
    description={Linguaggio di scripting orientato agli oggetti ed eventi comunemente utilizzato nella programmazione web lato client, anche se è stato recentemente esteso anche al lato server}
}

\newglossaryentry{oracle}
{
    name=Database Oracle,
    text=Database Oracle,
    sort=Database Oracle,
    description={È uno dei più noti \gls{dbms} relazionali, è scritto in linguaggio C. È prodotto dalla \emph{Oracle Corporation}}
}

\newglossaryentry{httpg}
{
    name=HTTP,
    text=HTTP,
    sort=HTTP,
    description={In telecomunicazioni, l'HyperText Transfer Protocol è un protocollo usato come sistema per la trasmissione di informazioni sul web}
}

\newglossaryentry{microservizio}
{
    name=Microservizio,
    text=microservizio,
    sort=Microservizio,
    description={Componente di un'architetettura a microservizi, variante dell'architettura orientata ai servizi, in cui l'applicazione è strutturata come un insieme di servizi scarsamente accoppiati}
}

\newglossaryentry{restg}
{
    name=REST,
    text=REST,
    sort=REST,
    description={Representational State Transfer è uno stile architetturale usato nei sistemi distribuiti che si riferisce a un sistema di trasmissione dati via \acrshort{http} senza \emph{layer} addizionali. I sistemi REST infatti non prevedono il concetto di sessione, essendo \textit{stateless}}
}

\newglossaryentry{spring-cloud}
{
    name=Spring Cloud,
    text=Spring Cloud,
    sort=Spring Cloud,
    description={\textit{Framework} del team Pivotal, sotto-progetto di \gls{spring}, che offre gli strumenti necessari agli sviluppatori per implementare alcuni dei \textit{pattern} più comuni per i sistemi distribuiti}
}

\newglossaryentry{typescript}
{
    name=TypeScript,
    text=TypeScript,
    sort=TypeScript,
    description={Linguaggio di programmazione \emph{open source} sviluppato da \emph{Microsoft}, super-set di \gls{javascript}, che estende la sintassi rendendolo un linguaggio tipizzato}
}

\newglossaryentry{deploy}
{
    name=Deploy,
    text=deploy,
    sort=Deploy,
    description={Nell'informatica, è un termine che si riferisce al \emph{deployment}, ossia la consegna o rilascio al cliente finale con eventuale installazione e messa in funzione di un sistema software}
}


\newglossaryentry{mocking}
{
    name=Mocking,
    text=mocking,
    sort=Mocking,
    description={Nell'informatica, i \emph{mock object} sono oggetti che simulano e riproducono il comportamento degli oggetti reali. Vengono usati per testare il comportamento di altri oggetti legati a un oggetto talvolta inaccessibile o non implementato}
}

\newglossaryentry{srp}
{
    name=Single Responsibility Principle,
    text=Single Responsibility Principle,
    sort=Single Responsibility Principle,
    description={Nella programmazione orientata agli oggetti questo principio afferma che ogni elemento di un sistema (variabile, classe, metodo) deve avere una e una sola responsabilità, e che tale responsabilità deve essere incapsulata dall'elemento stesso}
}

\newglossaryentry{design-pattern}
{
    name=Design Pattern,
    text=design pattern,
    sort=Design Pattern,
    description={Nell'ingegneria del software, un \emph{design pattern} è una soluzione progettuale generale ad un problema ricorrente. Si tratta, in pratica, di un modello astratto da applicare per la risoluzione di un problema che può presentarsi in diverse situazioni durante le fasi di progettazione e sviluppo di un sistema software}
}

\newglossaryentry{information-hiding}
{
    name=Information Hiding,
    text=information hiding,
    sort=Information Hiding,
    description={Principio della programmazione ad oggetti che si riferisce all'uso dell'incapsulamento per limitare l'accesso diretto agli elementi di un oggetto, per evitare danni collaterali nell'uso di dati inconsistenti con la logica prevista del programma}
}

\newglossaryentry{container}
{
    name=Container,
    text=container,
    sort=Container,
    description={Nell'informatica, un \emph{container} è un'istanza isolata nello spazio utente di un componente o sistema}
}

\newglossaryentry{macchina-virtuale}
{
    name=Macchina Virtuale,
    text=macchina virtuale,
    sort=Macchina Virtuale,
    description={In informatica, con macchina virtuale ci si riferisce a un software che crea un ambiente virtuale che emula tipicamente il comportamento di una macchina fisica tramite un processo di virtualizzazione, grazie all'assegnazione di risorse \emph{hardware}, quali \gls{ram}, porzioni di disco rigido, etc\dots }
}

\newglossaryentry{snippet}
{
    name=Snippet,
    text=snippet,
    sort=Snippet,
    description={Frammento di codice sorgente, generalmente distribuiti nel pubblico dominio}
}

\newglossaryentry{git}
{
    name=Git,
    text=Git,
    sort=Git,
    description={Sistema di controllo versione distribuito, utilizzabile da \gls{cli}, creato da Linus Torvalds nel 2005}
}

\newglossaryentry{repository}
{
    name=Repository,
    text=repository,
    sort=Repository,
    description={Ambiente di un sistema informatico in cui vengono gestiti metadati attraverso tabelle relazionali. Può essere implementato attraverso numerosi sistemi di gestione delle basi di dati}
}

\newglossaryentry{message-broker}
{
    name=Message Broker,
    text=message broker,
    sort=Message Broker,
    description={È un componente che traduce un messaggio dalla rappresentazione formale di protocollo di un messaggio del mittente a quello del ricevente}
}

\newglossaryentry{dependency-injection}
{
    name=Dependency Injection,
    text=dependency injection,
    sort=Dependency Injection,
    description={\glsdisp{design-pattern}{Design pattern}\gloss\ architetturale della programmazione a oggetti in cui un oggetto fornisce le dipendenze a un altro oggetto. Una ``dipendenza'' può essere vista come un oggetto che può essere usato, come un servizio}
}
 % database di termini
\makeglossaries


%**************************************************************
% Impostazioni di graphicx
%**************************************************************
\graphicspath{{immagini/}} % cartella dove sono riposte le immagini


%**************************************************************
% Impostazioni di hyperref
%**************************************************************
\hypersetup{
    %hyperfootnotes=false,
    %pdfpagelabels,
    %draft,	% = elimina tutti i link (utile per stampe in bianco e nero)
    colorlinks=true,
    linktocpage=true,
    pdfstartpage=1,
    pdfstartview=FitV,
    % decommenta la riga seguente per avere link in nero (per esempio per la stampa in bianco e nero)
    %colorlinks=false, linktocpage=false, pdfborder={0 0 0}, pdfstartpage=1, pdfstartview=FitV,
    breaklinks=true,
    pdfpagemode=UseNone,
    pageanchor=true,
    pdfpagemode=UseOutlines,
    plainpages=false,
    bookmarksnumbered,
    bookmarksopen=true,
    bookmarksopenlevel=1,
    hypertexnames=true,
    pdfhighlight=/O,
    %nesting=true,
    %frenchlinks,
    urlcolor=webbrown,
    linkcolor=RoyalBlue,
    citecolor=webgreen,
    %pagecolor=RoyalBlue,
    %urlcolor=Black, linkcolor=Black, citecolor=Black, %pagecolor=Black,
    pdftitle={\myTitle},
    pdfauthor={\textcopyright\ \myName, \myUni, \myFaculty},
    pdfsubject={},
    pdfkeywords={},
    pdfcreator={pdfLaTeX},
    pdfproducer={LaTeX}
}

%**************************************************************
% Impostazioni di itemize
%**************************************************************
\renewcommand{\labelitemi}{$\ast$}

%\renewcommand{\labelitemi}{$\bullet$}
%\renewcommand{\labelitemii}{$\cdot$}
%\renewcommand{\labelitemiii}{$\diamond$}
%\renewcommand{\labelitemiv}{$\ast$}


%**************************************************************
% Impostazioni di listings
%**************************************************************
\lstset{
    language=[LaTeX]Tex,%C++,
    keywordstyle=\color{RoyalBlue}, %\bfseries,
    basicstyle=\small\ttfamily,
    %identifierstyle=\color{NavyBlue},
    commentstyle=\color{Green}\ttfamily,
    stringstyle=\rmfamily,
    numbers=none, %left,%
    numberstyle=\scriptsize, %\tiny
    stepnumber=5,
    numbersep=8pt,
    showstringspaces=false,
    breaklines=true,
    frameround=ftff,
    frame=single
}


%**************************************************************
% Impostazioni di xcolor
%**************************************************************

\definecolor{webgreen}{rgb}{0,.5,0}
\definecolor{webbrown}{rgb}{.6,0,0}


%**************************************************************
% Altro
%**************************************************************

\newcommand{\omissis}{[\dots\negthinspace]} % produce [...]

% eccezioni all'algoritmo di sillabazione
\hyphenation
{
    ma-cro-istru-zio-ne
    gi-ral-din
}

\newcommand{\sectionname}{sezione}
\addto\captionsitalian{\renewcommand{\figurename}{Figura}
                       \renewcommand{\tablename}{Tabella}}

\newcommand{\glsfirstoccur}{\ap{{[g]}}}

\newcommand{\intro}[1]{\emph{\textsf{#1}}}

%**************************************************************
% Environment per ``rischi''
%**************************************************************
\newcounter{riskcounter}                % define a counter
\setcounter{riskcounter}{0}             % set the counter to some initial value

%%%% Parameters
% #1: Title
\newenvironment{risk}[1]{
    \refstepcounter{riskcounter}        % increment counter
    \par \noindent                      % start new paragraph
    \textbf{\arabic{riskcounter}. #1}   % display the title before the
                                        % content of the environment is displayed
}{
    \par\medskip
}

\newcommand{\riskname}{Rischio}

\newcommand{\riskdescription}[1]{\textbf{\\Descrizione:} #1.}

\newcommand{\risksolution}[1]{\textbf{\\Soluzione:} #1.}

%**************************************************************
% Environment per ``use case''
%**************************************************************
\newcounter{usecasecounter}             % define a counter
\setcounter{usecasecounter}{0}          % set the counter to some initial value

%%%% Parameters
% #1: ID
% #2: Nome
\newenvironment{usecase}[2]{
    \renewcommand{\theusecasecounter}{\usecasename #1}  % this is where the display of
                                                        % the counter is overwritten/modified
    \refstepcounter{usecasecounter}             % increment counter
    \vspace{10pt}
    \par \noindent                              % start new paragraph
    {\large \textbf{\usecasename #1: #2}}       % display the title before the
                                                % content of the environment is displayed
    \medskip
}{
    \medskip
}

\newcommand{\usecasename}{UC}

\newcommand{\usecaseactors}[1]{\textbf{\\Attori Principali:} #1. \vspace{4pt}}
\newcommand{\usecasepre}[1]{\textbf{\\Precondizioni:} #1. \vspace{4pt}}
\newcommand{\usecasedesc}[1]{\textbf{\\Descrizione:} #1. \vspace{4pt}}
\newcommand{\usecasepost}[1]{\textbf{\\Postcondizioni:} #1. \vspace{4pt}}
\newcommand{\usecasealt}[1]{\textbf{\\Scenario Alternativo:} #1. \vspace{4pt}}

%**************************************************************
% Environment per ``namespace description''
%**************************************************************

\newenvironment{namespacedesc}{
    \vspace{10pt}
    \par \noindent                              % start new paragraph
    \begin{description}
}{
    \end{description}
    \medskip
}

\newcommand{\classdesc}[2]{\item[\textbf{#1:}] #2}
                     % file con le impostazioni personali

\begin{document}

%**************************************************************
% Materiale iniziale
%**************************************************************

\frontmatter
% !TEX encoding = UTF-8
% !TEX TS-program = pdflatex
% !TEX root = ../tesi.tex

%**************************************************************
% Frontespizio
%**************************************************************
\begin{titlepage}

\begin{center}

\begin{LARGE}
\textbf{\myUni}\\
\end{LARGE}

\vspace{10pt}

\begin{Large}
\textsc{\myDepartment}\\
\end{Large}

\vspace{10pt}

\begin{large}
\textsc{\myFaculty}\\
\end{large}

\vspace{30pt}
\begin{figure}[htbp]
\begin{center}
\includegraphics[height=6cm]{immagini/logo-unipd.png}
\end{center}
\end{figure}
\vspace{30pt}

\begin{LARGE}
\begin{center}
\textbf{\myTitle}\\
\end{center}
\end{LARGE}

\vspace{10pt}

\begin{large}
\textsl{\myDegree}\\
\end{large}

\vspace{40pt}

\begin{large}
\begin{flushleft}
\textit{Relatore}\\
\vspace{5pt}
\profTitle \ \myProf
\end{flushleft}

\vspace{0pt}

\begin{flushright}
\textit{Laureando}\\
\vspace{5pt}
\myName
\end{flushright}
\end{large}

\vspace{40pt}

\line(1, 0){338} \\
\begin{normalsize}
\textsc{Anno Accademico \myAA}
\end{normalsize}

\end{center}
\end{titlepage}
% !TEX encoding = UTF-8
% !TEX TS-program = pdflatex
% !TEX root = ../tesi.tex

%**************************************************************
% Colophon
%**************************************************************
\clearpage
\phantomsection
\thispagestyle{empty}

\hfill

\vfill

\noindent\myName: \textit{\myTitle},
\myDegree,
\textcopyright\ \myTime.
\input{inizio-fine/dedica}
% !TEX encoding = UTF-8
% !TEX TS-program = pdflatex
% !TEX root = ../tesi.tex

%**************************************************************
% Sommario
%**************************************************************
\cleardoublepage
\phantomsection
\pdfbookmark{Sommario}{Sommario}
\begingroup
\let\clearpage\relax
\let\cleardoublepage\relax
\let\cleardoublepage\relax

\chapter*{Sommario}

Il presente documento descrive il lavoro svolto durante il periodo di stage, della durata di circa trecento ore, dal laureando \myName\ presso l'azienda \myCompany\
Gli obiettivi da raggiungere erano molteplici.

In primo luogo, era previsto lo studio delle tecnologie necessarie per poter approcciare l'implementazione del progetto di stage. Tra queste, sono state centrali
\textit{Java} e \textit{Spring}, con qualche accenno a tecnologie per il front-end, quali \textit{Typescript} e \textit{Angular}.

\noindent È stata dedicata inoltre una parte allo studio dell'architettura a microservizi e a \textit{Spring Cloud}.

In secondo luogo era richiesta l'implementazione del progetto di stage, che è consistito nello sviluppo del back-end di una \textit{web application} per gestire lo \textit{skill matrix}, utilizzato dall'azienda per raccogliere le competenze dei candidati ai colloqui lavorativi.

%\vfill
%
%\selectlanguage{english}
%\pdfbookmark{Abstract}{Abstract}
%\chapter*{Abstract}
%
%\selectlanguage{italian}

\endgroup

\vfill

% !TEX encoding = UTF-8
% !TEX TS-program = pdflatex
% !TEX root = ../tesi.tex

%**************************************************************
% Ringraziamenti
%**************************************************************
\cleardoublepage
\phantomsection
\pdfbookmark{Ringraziamenti}{ringraziamenti}

% \begin{flushright}{
% 	\slshape
% 	``Life is really simple, but we insist on making it complicated''} \\
% 	\medskip
%     --- Confucius
% \end{flushright}


\bigskip

\begingroup
\let\clearpage\relax
\let\cleardoublepage\relax
\let\cleardoublepage\relax

\chapter*{Ringraziamenti}

\noindent \textit{Innanzitutto, vorrei esprimere la mia gratitudine al Prof. \myProf, relatore della mia tesi, per l'aiuto e il sostegno fornitomi durante la stesura del lavoro.}
\bigskip

\noindent \textit{Ho desiderio di ringraziare il mio tutor aziendale, \fabio, per avermi seguito durante il periodo di stage e consigliato nei momenti critici.} % TODO
\bigskip

\noindent \textit{Desidero ringraziare con affetto i miei genitori e amici per il sostegno, il grande aiuto e per essermi stati vicini in ogni momento durante gli anni di studio.}
\bigskip


\noindent\textit{\myLocation, \myTime}
\hfill \myName

\endgroup

% !TEX encoding = UTF-8
% !TEX TS-program = pdflatex
% !TEX root = ../tesi.tex

%**************************************************************
% Indici
%**************************************************************
\cleardoublepage
\pdfbookmark{\contentsname}{tableofcontents}
\setcounter{tocdepth}{2}
\tableofcontents
%\markboth{\contentsname}{\contentsname} 
\clearpage

\begingroup 
    \let\clearpage\relax
    \let\cleardoublepage\relax
    \let\cleardoublepage\relax
    %*******************************************************
    % Elenco delle figure
    %*******************************************************    
    \phantomsection
    \pdfbookmark{\listfigurename}{lof}
    \listoffigures

    \vspace*{4ex}

    %*******************************************************
    % Elenco delle tabelle
    %*******************************************************
    \phantomsection
    \pdfbookmark{\listtablename}{lot}
    \listoftables
    
    \vspace*{4ex}
    
    \phantomsection
    \pdfbookmark{\listcodesname}{codes}
    \listofcodes
    
\endgroup

\cleardoublepage

\cleardoublepage

%**************************************************************
% Materiale principale
%**************************************************************

\mainmatter
% !TEX encoding = UTF-8
% !TEX TS-program = pdflatex
% !TEX root = ../tesi.tex

%**************************************************************
\chapter{Introduzione}
\label{cap:introduzione}
%**************************************************************

Introduzione al contesto applicativo.\bigskip

Sarà presentata nel presente capitolo l'azienda ospitante, Sync Lab, e la descrizione del progetto di stage.

\noindent Esempio di utilizzo di un termine nel glossario \par
\noindent \gls{api}.\bigskip

\noindent Esempio di citazione in linea \par
\noindent \cite{site:agile-manifesto}. \bigskip

\noindent Esempio di citazione nel pie' di pagina \par
\noindent citazione\footcite{womak:lean-thinking}

%**************************************************************

\section{Convenzioni tipografiche}

% \begin{description}
% \item[{\hyperref[cap:processi-metodologie]{Il secondo capitolo}}] descrive ...

% \item[{\hyperref[cap:descrizione-stage]{Il terzo capitolo}}] approfondisce ...

% \item[{\hyperref[cap:analisi-requisiti]{Il quarto capitolo}}] approfondisce ...

% \item[{\hyperref[cap:progettazione-codifica]{Il quinto capitolo}}] approfondisce ...

% \item[{\hyperref[cap:verifica-validazione]{Il sesto capitolo}}] approfondisce ...

% \item[{\hyperref[cap:conclusioni]{Nel settimo capitolo}}] descrive ...
% \end{description}

Riguardo la stesura del testo, relativamente al documento sono state adottate le seguenti convenzioni tipografiche:
\begin{itemize}
	\item Gli acronimi, le abbreviazioni e i termini ambigui o di uso non comune menzionati vengono definiti nel glossario, situato alla fine del presente documento;
	\item Per la prima occorrenza dei termini riportati nel glossario viene utilizzata la seguente nomenclatura: \gloss{parola}.
	\item I termini in lingua straniera o facenti parti del gergo tecnico sono evidenziati con il carattere \emph{corsivo}.
\end{itemize}

%**************************************************************

\section{L'azienda}

\begin{figure}[H]
	\centering
	\includegraphics[width=\textwidth/2]{immagini/logo-synclab.png}
	\caption{Logo Sync Lab}
\end{figure}

\subsection{Profilo aziendale}
\myCompany\ è una società di consulenza informatica fondata nel 2002 con sedi a Napoli, Roma, Milano e Padova.

Fin dai primi anni Sync Lab è rapidamente cresciuta nel mercato ICT, consolidando i rapporti con clienti e partner, raggiungendo un organico aziendale di oltre 200 risorse,
una solida base finanziaria e una diffusione sul territorio attraverso le sue quattro sedi.

L'organico aziendale è andato crescendo in modo continuo e rapido, in relazione
all'apertura delle varie sedi ed alla progressiva crescita delle stesse.

La grande attenzione alla gestione delle risorse umane ha fatto di Sync Lab un
riferimento in positivo per quanti volessero avviare o far evolvere in chiave professionale la propria carriera.

Il basso turn-over testimonia la voglia dei collaboratori di condividere il progetto comune, assumendo all'interno di esso ruoli e responsabilità che solo un processo
evolutivo così intenso può offrire.
I ricavi hanno avuto un incremento proporzionale alla crescita dell'azienda beneficiando dell'approccio adattativo e diversificato al mercato.

\subsection{Servizi}

\myCompany\ è un'azienda leader nella consulenza tecnologica, impegnata in un
processo continuo di identificazione e messa in opera di soluzioni per i clienti finalizzate
alla creazione di valore.

L'azienda supporta le esigenze di innovazione di tutte le
organizzazioni ed in ogni settore di mercato nell'ambito \textit{Information Technology}, con
servizi in ambito:
\begin{itemize}
	\item \textit{Business Consultancy};
	\item \textit{Project Financing};
	\item \textit{IT Consultancy};
\end{itemize}

Sync Lab ha come punti di forza la qualità dei servizi offerti (certificazioni ISO 9001,
ISO 14001, ISO 27001, OHSAS 18001) ed un'accurata gestione delle risorse umane.
L'approfondita conoscenza di processi e tecnologie, maturata in esperienze altamente
significative e qualificanti, fornisce l'\textit{expertise} e il \textit{know-how} necessari per gestire
progetti di elevata complessità, dominando l'intero ciclo di vita: Studio di fattibilità,
Progettazione, Implementazione, \textit{Governance} e \textit{Post Delivery}.

L'offerta di consulenza specialistica trova le punte di eccellenza nella progettazione di
architetture software avanzate, siano esse per applicativi di dominio, per sistemi di
supporto al business (BSS), per sistemi di integrazione (EAI/SOA) o per sistemi di
monitoraggio applicativo/territoriale.

Il laboratorio RD (Ricerca e Sviluppo) dell'azienda è sempre al passo con i nuovi
paradigmi tecnologici e di comunicazione, ad esempio \textit{Big Data}, \textit{Cloud Computing},
\textit{Internet of Things}, \textit{Mobile e Sicurezza IT}, per supportare i propri clienti nella creazione
ed integrazione di applicazioni, processi e dispositivi.

Le attività in ambito \textit{Educational} ed RD hanno permesso di acquisire una profonda
conoscenza degli strumenti di finanza agevolata fruendone direttamente ed interagendo
con enti di supporto ai progetti innovativi dei propri clienti. L'azienda, grazie alla rete
di relazioni a livello nazionale ed internazionale, ha ottenuto importanti finanziamenti
in progetti RD europei (FP7 e H2020).

\subsection{Settori d'impiego}

Sync Lab si sta sempre più specializzando in vari settori d'impiego: dal mondo \textit{banking}
all'\textit{assurance} con una nicchia importante nell'ambito sanità in cui vanta un prodotto
d'eccellenza per la gestione delle cliniche private.
L'azienda inoltre ha recentemente fondato una collegata \textbf{Sync Security} che si occupa
espressamente del mondo della \textit{cyber security} e sicurezza informatica in genere.

%**************************************************************

\section{Lo stage}

Lo stage si è svolto nella sede di Padova dell'azienda ospitante, Sync Lab, nell'arco delle otto settimane previste nel piano di lavoro:
\begin{itemize}
	\item inizio il 6 maggio 2019;
	\item fine il 28 giugno 2019.
\end{itemize}
Il progetto di stage consiste nell'implementazione di una \textit{web application} a microservizi per il trasferimento in una nuova piattaforma del servizio di \textit{skill matrix} usato da Sync Lab per organizzare le competenze delle persone che vanno in azienda per un colloquio.

Prima di ogni colloquio a Sync Lab, è richiesto al candidato di compilare un foglio di calcolo contenente le competenze di interesse
per l'azienda e il livello, valore che il candidato dovrà indicare.

L'idea è quella di un portale web in cui il candidato compilerà i livelli delle competenze, in alternativa al vecchio foglio di calcolo.
Per fare ciò, è stato scelto un approccio moderno: 
\begin{itemize}
	\item il \textit{front-end} utilizzerà \textit{framework} di ultima generazione (nel nostro caso, Angular);
	\item il \textit{back-end} utilizzerà un'architettura a microservizi, i quali saranno implementati con il framework Spring (Java), per questioni di uniformità e semplicità, con MongoDB come DBMS.
\end{itemize}

A svolgere tale progetto, ci saranno altri stagisti oltre al sottoscritto, a cui a ognuno sarà assegnato un equo carico di lavoro.

Per quanto riguarda il back-end, l'azienda possiede già un portale web usato per caricare i dati dei candidati. Il DBMS utilizzato per i dati di tale portale è Oracle SQL.
Parte del lavoro dello stage consisterà nello spostare e riadattare il codice del vecchio portale a quello del nuovo.

Il mio ruolo nel progetto consisterà prettamente nell'implementazione di alcuni microservizi back-end e nel riadattamento del vecchio portale.
%**************************************************************
             % Introduzione
% \input{capitoli/capitolo-2}             % Processi
% !TEX encoding = UTF-8
% !TEX TS-program = pdflatex
% !TEX root = ../tesi.tex

%**************************************************************
\chapter{Descrizione dello stage}\label{cap:descrizione-stage}
%**************************************************************

\intro{Breve introduzione al capitolo}

%**************************************************************
\section{Introduzione al progetto}

%**************************************************************
\section{Analisi preventiva dei rischi}

Durante la fase di analisi iniziale sono stati individuati alcuni possibili rischi a cui si potrà andare incontro.
Si è quindi proceduto a elaborare delle possibili soluzioni per far fronte a tali rischi.

\begin{risk}{Performance del portatile utilizzato}
    \riskdescription{le performance del portatile, essendo ormai datato, potrebbero risultare lente o non abbastanza buone da causare dei rallentamenti sulle
        attività lavorative e di studio durante il periodo di stage}
    \risksolution{cercare di limitare gli sprechi di memoria utilizzando applicazioni non necessarie, e coinvolgere il tutor aziendale per l'eventuale sostituzione
        del PC}
    \label{risk:hardware-simulator}
\end{risk}

%**************************************************************
\section{Requisiti e obiettivi}


%**************************************************************
\section{Pianificazione}
             % Kick-Off
\input{capitoli/capitolo-4}             % Concept Preview
\input{capitoli/capitolo-5}             % Product Prototype
\input{capitoli/capitolo-6}             % Product Design Freeze e SOP
\input{capitoli/capitolo-7}             % Conclusioni
\appendix
\input{capitoli/capitolo-A}             % Appendice A

%**************************************************************
% Materiale finale
%**************************************************************

\backmatter
% !TEX encoding = UTF-8
% !TEX TS-program = pdflatex
% !TEX root = ../tesi.tex

%**************************************************************
% Bibliografia
%**************************************************************

\cleardoublepage
\chapter{Bibliografia}

%\nocite{*}
% Stampa i riferimenti bibliografici
%\printbibliography[heading=subbibliography,title={Riferimenti bibliografici},type=book]

% Stampa i siti web consultati
\printbibliography[heading=subbibliography,title={Siti web consultati},type=online]


\glsaddall
\printglossaries
\end{document}
